% !TeX program = xelatex
% ============================================================
% Аннотация — документ КТ1 исследовательской (академической) ВКР
% ОП «Программная инженерия» ФКН НИУ ВШЭ.
%
% ЖАНР (Инструкция по подготовке документов к защите тем ВКР 2025/26 +
% ГОСТ 7.32-2017 §5.3): набор ПОЛУЖИРНЫХ ЛЕЙБЛОВ — структурных частей
% аннотации — со значением следом. НЕ ЕСПД: без штампа, без кода
% документа, без листа утверждения, без содержания и без нумерованных
% разделов. Требования КТ1: аннотация СО СПИСКОМ ИСТОЧНИКОВ — не менее
% 7–8 страниц; источников — не менее 5, по ГОСТ Р 7.0.5-2008 (алфавит,
% латиница → кириллица); ТИТУЛЬНЫЙ ЛИСТ ПОДПИСЫВАЮТ студент И
% руководитель (подписи размещаются в разделе «Подписи»).
%
% Как пользоваться: замените жёлтые \hseFill{…} своим текстом. Если по
% какой-то структурной части сведений нет — часть опускается, при этом
% ПОСЛЕДОВАТЕЛЬНОСТЬ остальных сохраняется (инструкция §3). Данные о
% себе, руководителе и годе задаются в настройках проекта (макросы \hse…).
% ============================================================
\documentclass[12pt,a4paper]{extarticle}
% !TeX program = xelatex
% ============================================================
% common/annotation_preamble.tex — преамбула Аннотации (документ КТ1
% исследовательской ВКР ОП «Программная инженерия»).
%
% ЭТО НЕ ЕСПД-документ: без штампа/рамки, без кода документа, без листа
% утверждения и без содержания. Оформление по ГОСТ 7.32-2017: поля
% 30/15/20/20, Times New Roman 12 пт, интервал 1,5, обычный номер
% страницы снизу. Данные — из common/meta (\hse-макросы). Ссылки [n]
% печатаются В СТРОКУ, список источников — по ГОСТ Р 7.0.5-2008.
% ============================================================

% ── Шрифты (автогенерируются из настроек проекта) ───────────
\newcommand{\hseDefaultMono}{Courier New}
\input{../common/fonts}
\usepackage[english,russian]{babel}

% ── Геометрия по ГОСТ 7.32-2017 §6.1.1 (без места под штамп) ─
\usepackage[a4paper,left=30mm,right=15mm,top=20mm,bottom=20mm]{geometry}
\usepackage{setspace}
\onehalfspacing
\usepackage{indentfirst}
\setlength{\parindent}{1.25cm}
\setlength{\parskip}{6pt}

% ── Типографика и таблицы ───────────────────────────────────
\usepackage{microtype}
\usepackage{csquotes}
\usepackage{enumitem}
\setlist[itemize,1]{label=--}
\usepackage{graphicx}
\usepackage[table]{xcolor}
\usepackage{array}
\usepackage{tabularx}
\usepackage{booktabs}
\usepackage{float}

% ── URL и переносы ──────────────────────────────────────────
\usepackage{url}
\usepackage{xurl}
\urlstyle{same}
\usepackage{hyphenat}

% ── Цитирование: ссылки [n] В СТРОКУ (как в реальных аннотациях),
%    список источников — ГОСТ Р 7.0.5-2008. НЕ переопределяем \@cite
%    в надстрочный вид (в отличие от ЕСПД-преамбулы). ────────────
\usepackage{cite}
\makeatletter
\renewcommand\@biblabel[1]{#1.}
\makeatother

\usepackage{lastpage}
\usepackage{hyperref}
\hypersetup{
    unicode=true,
    breaklinks=true,
    colorlinks=true,
    linkcolor=black,
    urlcolor=blue,
    citecolor=blue,
    pdfborder={0 0 0}
}

% ── Данные проекта (\hse…), затем «человеческие» имена (commands) ──
% meta даёт \hseFill/\hseOptional и все \hse-макросы; commands —
% \signdateblank, \hseExecutorsBlock-зависимости (\authorsignatureline,
% \studentgroup, \studentname) для блока «Выполнил» в титуле.
\input{../common/meta}
\input{../common/commands}

% Без штампа: обычный номер страницы снизу по центру (ГОСТ 7.32 §6.3.1).
\pagestyle{plain}

\begin{document}
\sloppy
% ============================================================
% common/annotation_title.tex — титульный лист Аннотации ВКР ОП ПИ
% по официальной форме «Титул Аннотация защита темы ВКР ПИ».
%
% Лёгкий титул БЕЗ штампа/рамки/листа утверждения: УДК, две колонки
% СОГЛАСОВАНО (руководитель проекта) + УТВЕРЖДАЮ (академический
% руководитель ОП), затем «Аннотация выпускной квалификационной
% работы», тема, направление и блок «Выполнил» (в командной работе —
% состав исполнителей через \hseExecutorsBlock).
% ============================================================
\begin{titlepage}
\thispagestyle{empty}
\singlespacing
\footnotesize
\setlength{\parindent}{0pt}
\setlength{\parskip}{0pt}

\begin{center}
\textbf{ПРАВИТЕЛЬСТВО РОССИЙСКОЙ ФЕДЕРАЦИИ}\\
\textbf{ФЕДЕРАЛЬНОЕ ГОСУДАРСТВЕННОЕ АВТОНОМНОЕ}\\
\textbf{ОБРАЗОВАТЕЛЬНОЕ УЧРЕЖДЕНИЕ ВЫСШЕГО ОБРАЗОВАНИЯ}\\
\textbf{НАЦИОНАЛЬНЫЙ ИССЛЕДОВАТЕЛЬСКИЙ УНИВЕРСИТЕТ}\\
\textbf{«ВЫСШАЯ ШКОЛА ЭКОНОМИКИ»}

\vspace{0.35cm}
\hseFaculty\\
Образовательная программа «\hseSpecialization»
\end{center}

\vspace{0.25cm}
\noindent
УДК~\hseUDC

\vspace{0.4cm}
\noindent
\begin{minipage}[t]{0.46\textwidth}
\raggedright
\textbf{СОГЛАСОВАНО}\\
Руководитель проекта

\vspace{0.15cm}
\strut\hseSupervisorTitle\\[-0.05cm]
\rule{6cm}{0.4pt}\\[-0.05cm]
{\scriptsize Должность, место работы, ученая степень}

\vspace{0.12cm}
\strut\hseSupervisorName\\[-0.05cm]
\rule{6cm}{0.4pt}\\[-0.05cm]
{\scriptsize Подпись, И.О. Фамилия}

\vspace{0.1cm}
\signdateblank
\end{minipage}
\hfill
\begin{minipage}[t]{0.46\textwidth}
\raggedright
\textbf{УТВЕРЖДАЮ}\\
\strut\hseAcademicSupervisorTitle

\vspace{0.35cm}
\mbox{\rule{6cm}{0.4pt}}\\[-0.05cm]
\hseAcademicSupervisorName

\vspace{0.12cm}
«\rule{0.8cm}{0.4pt}» \rule{2.5cm}{0.4pt} \hseYear~г.
\end{minipage}

\vspace{0.95cm}
\begin{center}
\textbf{\normalsize Аннотация}

\vspace{0.12cm}
\textbf{\normalsize выпускной квалификационной работы}

\vspace{0.35cm}
на тему
\underline{\makebox[0.72\textwidth][c]{\hseProjectName}}

\vspace{0.3cm}
по направлению подготовки бакалавров \hseProgramCode\ «\hseSpecialization»
\end{center}

\vspace{0.95cm}
\noindent\hfill
\begin{minipage}[t]{0.52\textwidth}
\raggedright
% Командная работа (shared-база) — состав исполнителей; solo/личная — один.
\ifhseSharedDoc
\hseExecutorsBlock
\else
Выполнил\\
студент группы \hseGroup\\
образовательной программы\\
\hseProgramCode\ «\hseSpecialization»

\vspace{0.25cm}
\strut\hseAuthorName\\[-0.05cm]
\rule{6cm}{0.4pt}\\[-0.05cm]
{\scriptsize И.О. Фамилия}

\vspace{0.15cm}
\rule{3.6cm}{0.4pt}\quad\rule{2.2cm}{0.4pt}\\[-0.05cm]
{\scriptsize Подпись, Дата}
\fi
\end{minipage}

\vfill
\begin{center}
\textbf{\hseCity{} \hseYear}
\end{center}
\end{titlepage}
\clearpage



\begin{center}
  \textbf{\large Аннотация}
\end{center}

\vspace{0.4em}

% ── Ключевые слова (ГОСТ 7.32-2017 §5.3.2.1 / §6.12.2): 5–15 слов или
%    словосочетаний, ПРОПИСНЫМИ буквами, в именительном падеже, в строку
%    через запятую, без точки в конце. ──
\noindent\textbf{Ключевые слова:} \hseFill{5–15 СЛОВ ПРОПИСНЫМИ В ИМЕНИТЕЛЬНОМ ПАДЕЖЕ ЧЕРЕЗ ЗАПЯТУЮ БЕЗ ТОЧКИ В КОНЦЕ}

\medskip
\noindent\textbf{Оценка современного состояния решаемой проблемы.} Рассматриваемая проблема состоит в~\hseFill{кратко сформулируйте научную проблему}. Сегодня для её решения применяются \hseFill{перечислите существующие подходы, методы, работы} \cite{example-en}. Их основные ограничения: \hseFill{ограничение 1}; \hseFill{ограничение 2}~--- что и~определяет направление настоящего исследования.

\medskip
\noindent\textbf{Основание и исходные данные для исследования.} Основанием для выполнения работы является учебный план ОП «\hseSpecialization» и~тема ВКР, утверждаемая приказом ФКН НИУ ВШЭ. Исходные данные: \hseFill{имеющийся задел, данные, публикации, программные средства, от которых отталкивается работа}.

\medskip
\noindent\textbf{Актуальность и новизна темы исследования.} Актуальность обусловлена \hseFill{практическая и научная значимость задачи}. Научная новизна работы заключается в~\hseFill{что именно предлагается впервые}.

\medskip
\noindent\textbf{Объект исследования:} \hseFill{что изучается в целом}.

\smallskip
\noindent\textbf{Предмет исследования:} \hseFill{конкретная сторона объекта, на которой сфокусирована работа}.

\smallskip
\noindent\textbf{Цель исследования:} \hseFill{одна формулировка достигаемого научного результата}.

\medskip
\noindent\textbf{Задачи исследования:}
\begin{enumerate}[label=\arabic*), leftmargin=1.3cm, topsep=2pt, itemsep=2pt]
  \item \hseFill{провести обзор существующих подходов};
  \item \hseFill{предложить и формализовать метод (модель, алгоритм)};
  \item \hseFill{реализовать программное средство для экспериментов};
  \item \hseFill{провести вычислительный эксперимент и оценить результаты};
  \item \hseFill{сделать выводы о применимости}.
\end{enumerate}

\medskip
\noindent\textbf{Методы исследования:} \hseFill{перечислите методы: аналитический обзор, формальное моделирование, вычислительный эксперимент, статистический анализ и т.\,п.}

\medskip
\noindent\textbf{Достоверность научных результатов} подтверждается \hseFill{воспроизводимостью экспериментов, сравнением с существующими подходами, статистической значимостью различий}.

\medskip
\noindent\textbf{Ожидаемые результаты и их новизна:}
\begin{enumerate}[label=\arabic*), leftmargin=1.3cm, topsep=2pt, itemsep=2pt]
  \item \hseFill{метод / модель / алгоритм};
  \item \hseFill{программная реализация для проведения экспериментов};
  \item \hseFill{экспериментальная оценка и сравнение с базовыми методами};
  \item \hseFill{выводы о границах применимости}.
\end{enumerate}

\medskip
\noindent\textbf{Область применения результатов:} \hseFill{укажите область применения}.

\medskip
\noindent\textbf{Предполагаемое внедрение результатов работы.} \hseFill{где и в каком виде планируется внедрить результаты: продукт, сервис, публикация, открытый код}.

\medskip
\noindent\textbf{Экономическая эффективность или значимость работы.} Результаты работы могут быть использованы \hseFill{где и кем применимы полученные результаты; в чём практический или экономический выигрыш}.

\medskip
\noindent\textbf{Прогнозные предположения о развитии объекта исследования.} \hseFill{как может развиваться объект исследования и какое продолжение возможно у работы}.

\bigskip
% ── Список источников по ГОСТ Р 7.0.5-2008 (методичка ФКН/ПИ): единый
%    алфавитный нумерованный список, сначала латиница, затем кириллица;
%    разделитель областей «. –»; для сетевых ресурсов — «[Электронный
%    ресурс] … Режим доступа: URL (дата обращения: ДД.ММ.ГГ)». На каждый
%    источник должна быть ссылка \cite{…} в тексте.
%    Инструкция КТ1: источников — НЕ МЕНЕЕ ПЯТИ. ──
\noindent\textbf{Список использованных источников}
\vspace{-0.3em}
\renewcommand{\refname}{}
\begin{thebibliography}{99}
  \bibitem{example-en}
  Author, A.~B. Title of the paper / A.~B.~Author, C.~D.~Coauthor // Proceedings of the Conference.~--- New York, NY, USA: ACM.~--- 2024.~--- С.~10--20.

  \bibitem{example-web}
  Название ресурса [Электронный ресурс].~--- Режим доступа: \url{https://example.org} (дата обращения: 01.01.2026).
\end{thebibliography}


\end{document}
